\section{原则二：工具的命名空间}

当 Agent 同时访问数十个 MCP Server 和数百个工具时，功能重叠或用途模糊的工具会导致 Agent 困惑，不知道该使用哪一个。

\subsection{Namespacing 策略}

\textbf{命名空间}（Namespacing）通过将相关工具归类到共同前缀下来划定边界，MCP 客户端有时会默认执行此操作。

\begin{table}[H]
\centering
\caption{命名空间的两个维度}
\begin{tabular}{p{0.2\textwidth}p{0.3\textwidth}p{0.35\textwidth}}
\toprule
\textbf{维度} & \textbf{示例} & \textbf{效果} \\
\midrule
按服务分组 & \texttt{asana\_search}, \texttt{jira\_search} & 区分不同平台的同类操作 \\
\midrule
按资源分组 & \texttt{asana\_projects\_search}, \texttt{asana\_users\_search} & 区分同一平台的不同资源 \\
\bottomrule
\end{tabular}
\end{table}

\begin{knowledgebox}{前缀式 vs 后缀式命名}
Anthropic 发现，选择前缀式（\texttt{asana\_search}）还是后缀式（\texttt{search\_asana}）命名对评估结果有\textbf{非平凡的影响}（non-trivial effects）。效果因 LLM 而异，建议根据自己的评估结果来选择命名方案。
\end{knowledgebox}

\subsection{Namespacing 的深层价值}

命名空间带来的好处不仅仅是组织性，更在于它可以：

\begin{enumerate}[leftmargin=1.5em]
    \item \textbf{减少加载到 Agent 上下文中的工具数量和描述}：通过名称反映任务的自然分工
    \item \textbf{将 Agent 的认知计算卸载到 Tool 调用本身}：减少 Agent 在上下文中的推理负担
    \item \textbf{降低 Agent 犯错的整体风险}：错误包括调用错误的工具、使用错误的参数、调用工具太少、或错误处理工具响应
\end{enumerate}

\subsection{本章小结}

合理的命名空间策略能帮助 Agent 在大量工具中快速定位正确选择。按服务和资源两个维度组织工具名称，并通过评估来选择最优的命名方案。


